% figs.tex —— 由 ch8/make_figs.py 產生，請勿手動修改。
% 重新產生：python ch8/make_figs.py（第 8 章，動畫 W801–W804）

\newcommand{\EightFigChi}{%
\begin{tikzpicture}[x=2.1cm,y=0.68cm,>=Stealth,font=\scriptsize,baseline=(current bounding box.center)]
\draw[gray!25] (-3,-5) grid (-1,0);
\draw[->] (-3,-5) -- (-0.8,-5) node[right] {$q$};
\draw[->] (-3,-5) -- (-3,0.35) node[above] {$1-\chi(q)$};
\foreach \x in {-3,-2,-1} {\node[below] at (\x,-5) {$10^{\x}$};}
\foreach \y in {-5,-4,-3,-2,-1,0} {\node[left] at (-3,\y) {$10^{\y}$};}
\draw[cBar,thin] plot coordinates {(-3.000,-3.105) (-2.750,-2.855) (-2.500,-2.605) (-2.250,-2.355) (-2.000,-2.106) (-1.750,-1.857) (-1.500,-1.608) (-1.250,-1.360) (-1.000,-1.114)};
\fill[cBar] (-3.000,-3.105) circle (1.5pt);
\fill[cBar] (-2.750,-2.855) circle (1.5pt);
\fill[cBar] (-2.500,-2.605) circle (1.5pt);
\fill[cBar] (-2.250,-2.355) circle (1.5pt);
\fill[cBar] (-2.000,-2.106) circle (1.5pt);
\fill[cBar] (-1.750,-1.857) circle (1.5pt);
\fill[cBar] (-1.500,-1.608) circle (1.5pt);
\fill[cBar] (-1.250,-1.360) circle (1.5pt);
\fill[cBar] (-1.000,-1.114) circle (1.5pt);
\fill[cBar] (-2.88,-0.15) circle (1.5pt); \node[cBar,anchor=west] at (-2.84,-0.15) {$\delta=1$：斜率 1.00};
\draw[cWalk,thin] plot coordinates {(-3.000,-4.506) (-2.750,-4.133) (-2.500,-3.761) (-2.250,-3.390) (-2.000,-3.021) (-1.750,-2.654) (-1.500,-2.289) (-1.250,-1.929) (-1.000,-1.573)};
\fill[cWalk] (-3.000,-4.506) circle (1.5pt);
\fill[cWalk] (-2.750,-4.133) circle (1.5pt);
\fill[cWalk] (-2.500,-3.761) circle (1.5pt);
\fill[cWalk] (-2.250,-3.390) circle (1.5pt);
\fill[cWalk] (-2.000,-3.021) circle (1.5pt);
\fill[cWalk] (-1.750,-2.654) circle (1.5pt);
\fill[cWalk] (-1.500,-2.289) circle (1.5pt);
\fill[cWalk] (-1.250,-1.929) circle (1.5pt);
\fill[cWalk] (-1.000,-1.573) circle (1.5pt);
\fill[cWalk] (-2.88,-0.6) circle (1.5pt); \node[cWalk,anchor=west] at (-2.84,-0.6) {$\delta=1/2$：斜率 1.49};
\draw[cExact,densely dashed,thick] plot coordinates {(-2.100,-4.978) (-2.000,-4.778) (-1.900,-4.578) (-1.800,-4.378) (-1.700,-4.178) (-1.600,-3.978) (-1.500,-3.778) (-1.400,-3.578) (-1.300,-3.378) (-1.200,-3.178) (-1.100,-2.978) (-1.000,-2.778)};
\draw[cExact,densely dashed,thick] (-2.94,-1.05) -- (-2.82,-1.05) node[anchor=west] {步長 1：斜率 2};
\end{tikzpicture}}

\newcommand{\EightTabWidth}{%
\begin{tabular}{@{}rrrrr@{}}
\toprule
$N$ & 長尾巴 & $\div N$ & 步長 1 & $\div\sqrt N$\\
\midrule
10 & 18 & 1.77 & 2.86 & 0.90\\
100 & 172 & 1.72 & 8.98 & 0.90\\
1000 & 1822 & 1.82 & 28.06 & 0.89\\
理論 & & 1.7785 & & 0.8881\\
斜率 & 1.01 & & 0.50 & \\
\bottomrule
\end{tabular}}

\newcommand{\EightTabTail}{%
\begin{tabular}{@{}lcccc@{}}
\toprule
$r$ & 10 & 30 & 100 & 斜率\\
\midrule
$\delta=1$ & 0.9803 & 0.9978 & 0.9998 & $-4.00$\\
$\delta=\frac12$ & 1.1755 & 1.0316 & 1.0051 & $-4.51$\\
\bottomrule
\end{tabular}}

\newcommand{\EightFigRawG}{%
\begin{tikzpicture}[x=1.6cm,y=0.068cm,>=Stealth,font=\scriptsize,baseline=(current bounding box.center)]
\draw[->] (-1.1,0) -- (1.2,0) node[right] {$q$};
\draw[->] (0,0) -- (0,55) node[above] {$g(z;q)$};
\foreach \x in {-1,1} {\draw (\x,0) -- (\x,-1.2); \node[below] at (\x,-1) {$\x$};}
\foreach \y in {10,30,50} {\draw (0,\y) -- (-0.04,\y); \node[left] at (-0.04,\y) {$\y$};}
\draw[cBar!40,thick] plot coordinates {(-1.000,1.761) (-0.975,1.815) (-0.950,1.870) (-0.925,1.929) (-0.900,1.989) (-0.875,2.053) (-0.850,2.119) (-0.825,2.187) (-0.800,2.259) (-0.775,2.334) (-0.750,2.412) (-0.725,2.493) (-0.700,2.576) (-0.675,2.664) (-0.650,2.754) (-0.625,2.847) (-0.600,2.944) (-0.575,3.043) (-0.550,3.145) (-0.525,3.249) (-0.500,3.356) (-0.475,3.465) (-0.450,3.576) (-0.425,3.688) (-0.400,3.800) (-0.375,3.912) (-0.350,4.024) (-0.325,4.134) (-0.300,4.242) (-0.275,4.347) (-0.250,4.447) (-0.225,4.542) (-0.200,4.631) (-0.175,4.712) (-0.150,4.785) (-0.125,4.849) (-0.100,4.902) (-0.075,4.944) (-0.050,4.975) (-0.025,4.994) (0.000,5.000) (0.025,4.994) (0.050,4.975) (0.075,4.944) (0.100,4.902) (0.125,4.849) (0.150,4.785) (0.175,4.712) (0.200,4.631) (0.225,4.542) (0.250,4.447) (0.275,4.347) (0.300,4.242) (0.325,4.134) (0.350,4.024) (0.375,3.912) (0.400,3.800) (0.425,3.688) (0.450,3.576) (0.475,3.465) (0.500,3.356) (0.525,3.249) (0.550,3.145) (0.575,3.043) (0.600,2.944) (0.625,2.847) (0.650,2.754) (0.675,2.664) (0.700,2.576) (0.725,2.493) (0.750,2.412) (0.775,2.334) (0.800,2.259) (0.825,2.187) (0.850,2.119) (0.875,2.053) (0.900,1.989) (0.925,1.929) (0.950,1.870) (0.975,1.815) (1.000,1.761)};
\draw[cBar!70,thick] plot coordinates {(-1.000,1.947) (-0.975,2.020) (-0.950,2.099) (-0.925,2.182) (-0.900,2.270) (-0.875,2.363) (-0.850,2.463) (-0.825,2.569) (-0.800,2.681) (-0.775,2.801) (-0.750,2.928) (-0.725,3.064) (-0.700,3.209) (-0.675,3.363) (-0.650,3.527) (-0.625,3.702) (-0.600,3.888) (-0.575,4.086) (-0.550,4.297) (-0.525,4.521) (-0.500,4.758) (-0.475,5.009) (-0.450,5.274) (-0.425,5.554) (-0.400,5.846) (-0.375,6.152) (-0.350,6.470) (-0.325,6.797) (-0.300,7.133) (-0.275,7.473) (-0.250,7.814) (-0.225,8.151) (-0.200,8.479) (-0.175,8.792) (-0.150,9.082) (-0.125,9.344) (-0.100,9.570) (-0.075,9.753) (-0.050,9.889) (-0.025,9.972) (0.000,10.000) (0.025,9.972) (0.050,9.889) (0.075,9.753) (0.100,9.570) (0.125,9.344) (0.150,9.082) (0.175,8.792) (0.200,8.479) (0.225,8.151) (0.250,7.814) (0.275,7.473) (0.300,7.133) (0.325,6.797) (0.350,6.470) (0.375,6.152) (0.400,5.846) (0.425,5.554) (0.450,5.274) (0.475,5.009) (0.500,4.758) (0.525,4.521) (0.550,4.297) (0.575,4.086) (0.600,3.888) (0.625,3.702) (0.650,3.527) (0.675,3.363) (0.700,3.209) (0.725,3.064) (0.750,2.928) (0.775,2.801) (0.800,2.681) (0.825,2.569) (0.850,2.463) (0.875,2.363) (0.900,2.270) (0.925,2.182) (0.950,2.099) (0.975,2.020) (1.000,1.947)};
\draw[cBar,thick] plot coordinates {(-1.000,2.125) (-0.975,2.222) (-0.950,2.326) (-0.925,2.438) (-0.900,2.559) (-0.875,2.689) (-0.850,2.831) (-0.825,2.985) (-0.800,3.152) (-0.775,3.335) (-0.750,3.534) (-0.725,3.753) (-0.700,3.993) (-0.675,4.257) (-0.650,4.549) (-0.625,4.872) (-0.600,5.231) (-0.575,5.631) (-0.550,6.078) (-0.525,6.580) (-0.500,7.144) (-0.475,7.783) (-0.450,8.506) (-0.425,9.330) (-0.400,10.271) (-0.375,11.350) (-0.350,12.592) (-0.325,14.025) (-0.300,15.681) (-0.275,17.598) (-0.250,19.815) (-0.225,22.370) (-0.200,25.294) (-0.175,28.598) (-0.150,32.254) (-0.125,36.171) (-0.100,40.167) (-0.075,43.946) (-0.050,47.115) (-0.025,49.246) (0.000,50.000) (0.025,49.246) (0.050,47.115) (0.075,43.946) (0.100,40.167) (0.125,36.171) (0.150,32.254) (0.175,28.598) (0.200,25.294) (0.225,22.370) (0.250,19.815) (0.275,17.598) (0.300,15.681) (0.325,14.025) (0.350,12.592) (0.375,11.350) (0.400,10.271) (0.425,9.330) (0.450,8.506) (0.475,7.783) (0.500,7.144) (0.525,6.580) (0.550,6.078) (0.575,5.631) (0.600,5.231) (0.625,4.872) (0.650,4.549) (0.675,4.257) (0.700,3.993) (0.725,3.753) (0.750,3.534) (0.775,3.335) (0.800,3.152) (0.825,2.985) (0.850,2.831) (0.875,2.689) (0.900,2.559) (0.925,2.438) (0.950,2.326) (0.975,2.222) (1.000,2.125)};
\node[cBar,anchor=west] at (0.06,50) {$z=0.49$};
\node[cBar!70,anchor=west] at (0.42,12.5) {$0.45$};
\node[cBar!40,anchor=west] at (0.72,7.5) {$0.4$};
\end{tikzpicture}}

\newcommand{\EightFigCollapse}{%
\begin{tikzpicture}[x=0.42cm,y=6.8cm,>=Stealth,font=\scriptsize,baseline=(current bounding box.center)]
\draw[->] (-4.3,0) -- (4.6,0) node[right] {$x$};
\draw[->] (0,0) -- (0,0.58) node[above] {$(z_c-z)\,g$};
\foreach \x in {-4,-2,2,4} {\draw (\x,0) -- (\x,-0.012); \node[below] at (\x,-0.01) {$\x$};}
\foreach \y in {0.1,0.25,0.5} {\draw (0,\y) -- (-0.12,\y); \node[left] at (-0.12,\y) {$\y$};}
\draw[cBar!40,thick] plot coordinates {(-4.000,0.1317) (-3.900,0.1360) (-3.800,0.1406) (-3.700,0.1454) (-3.600,0.1504) (-3.500,0.1558) (-3.400,0.1614) (-3.300,0.1674) (-3.200,0.1737) (-3.100,0.1803) (-3.000,0.1873) (-2.900,0.1948) (-2.800,0.2026) (-2.700,0.2108) (-2.600,0.2195) (-2.500,0.2287) (-2.400,0.2384) (-2.300,0.2485) (-2.200,0.2591) (-2.100,0.2703) (-2.000,0.2819) (-1.900,0.2940) (-1.800,0.3066) (-1.700,0.3196) (-1.600,0.3331) (-1.500,0.3468) (-1.400,0.3608) (-1.300,0.3750) (-1.200,0.3892) (-1.100,0.4034) (-1.000,0.4172) (-0.900,0.4307) (-0.800,0.4435) (-0.700,0.4555) (-0.600,0.4665) (-0.500,0.4762) (-0.400,0.4845) (-0.300,0.4912) (-0.200,0.4960) (-0.100,0.4990) (0.000,0.5000) (0.100,0.4990) (0.200,0.4960) (0.300,0.4912) (0.400,0.4845) (0.500,0.4762) (0.600,0.4665) (0.700,0.4555) (0.800,0.4435) (0.900,0.4307) (1.000,0.4172) (1.100,0.4034) (1.200,0.3892) (1.300,0.3750) (1.400,0.3608) (1.500,0.3468) (1.600,0.3331) (1.700,0.3196) (1.800,0.3066) (1.900,0.2940) (2.000,0.2819) (2.100,0.2703) (2.200,0.2591) (2.300,0.2485) (2.400,0.2384) (2.500,0.2287) (2.600,0.2195) (2.700,0.2108) (2.800,0.2026) (2.900,0.1948) (3.000,0.1873) (3.100,0.1803) (3.200,0.1737) (3.300,0.1674) (3.400,0.1614) (3.500,0.1558) (3.600,0.1504) (3.700,0.1454) (3.800,0.1406) (3.900,0.1360) (4.000,0.1317)};
\draw[cBar!70,thick] plot coordinates {(-4.000,0.1145) (-3.900,0.1187) (-3.800,0.1232) (-3.700,0.1279) (-3.600,0.1329) (-3.500,0.1382) (-3.400,0.1438) (-3.300,0.1497) (-3.200,0.1559) (-3.100,0.1625) (-3.000,0.1695) (-2.900,0.1769) (-2.800,0.1847) (-2.700,0.1930) (-2.600,0.2017) (-2.500,0.2110) (-2.400,0.2207) (-2.300,0.2310) (-2.200,0.2419) (-2.100,0.2533) (-2.000,0.2652) (-1.900,0.2778) (-1.800,0.2908) (-1.700,0.3044) (-1.600,0.3185) (-1.500,0.3330) (-1.400,0.3478) (-1.300,0.3630) (-1.200,0.3782) (-1.100,0.3934) (-1.000,0.4085) (-0.900,0.4231) (-0.800,0.4372) (-0.700,0.4504) (-0.600,0.4626) (-0.500,0.4734) (-0.400,0.4826) (-0.300,0.4901) (-0.200,0.4955) (-0.100,0.4989) (0.000,0.5000) (0.100,0.4989) (0.200,0.4955) (0.300,0.4901) (0.400,0.4826) (0.500,0.4734) (0.600,0.4626) (0.700,0.4504) (0.800,0.4372) (0.900,0.4231) (1.000,0.4085) (1.100,0.3934) (1.200,0.3782) (1.300,0.3630) (1.400,0.3478) (1.500,0.3330) (1.600,0.3185) (1.700,0.3044) (1.800,0.2908) (1.900,0.2778) (2.000,0.2652) (2.100,0.2533) (2.200,0.2419) (2.300,0.2310) (2.400,0.2207) (2.500,0.2110) (2.600,0.2017) (2.700,0.1930) (2.800,0.1847) (2.900,0.1769) (3.000,0.1695) (3.100,0.1625) (3.200,0.1559) (3.300,0.1497) (3.400,0.1438) (3.500,0.1382) (3.600,0.1329) (3.700,0.1279) (3.800,0.1232) (3.900,0.1187) (4.000,0.1145)};
\draw[cBar,thick] plot coordinates {(-4.000,0.1027) (-3.900,0.1069) (-3.800,0.1112) (-3.700,0.1158) (-3.600,0.1207) (-3.500,0.1259) (-3.400,0.1314) (-3.300,0.1372) (-3.200,0.1434) (-3.100,0.1499) (-3.000,0.1568) (-2.900,0.1641) (-2.800,0.1719) (-2.700,0.1802) (-2.600,0.1889) (-2.500,0.1982) (-2.400,0.2079) (-2.300,0.2183) (-2.200,0.2293) (-2.100,0.2408) (-2.000,0.2529) (-1.900,0.2657) (-1.800,0.2791) (-1.700,0.2930) (-1.600,0.3075) (-1.500,0.3225) (-1.400,0.3380) (-1.300,0.3537) (-1.200,0.3697) (-1.100,0.3858) (-1.000,0.4017) (-0.900,0.4173) (-0.800,0.4323) (-0.700,0.4464) (-0.600,0.4595) (-0.500,0.4711) (-0.400,0.4811) (-0.300,0.4892) (-0.200,0.4951) (-0.100,0.4988) (0.000,0.5000) (0.100,0.4988) (0.200,0.4951) (0.300,0.4892) (0.400,0.4811) (0.500,0.4711) (0.600,0.4595) (0.700,0.4464) (0.800,0.4323) (0.900,0.4173) (1.000,0.4017) (1.100,0.3858) (1.200,0.3697) (1.300,0.3537) (1.400,0.3380) (1.500,0.3225) (1.600,0.3075) (1.700,0.2930) (1.800,0.2791) (1.900,0.2657) (2.000,0.2529) (2.100,0.2408) (2.200,0.2293) (2.300,0.2183) (2.400,0.2079) (2.500,0.1982) (2.600,0.1889) (2.700,0.1802) (2.800,0.1719) (2.900,0.1641) (3.000,0.1568) (3.100,0.1499) (3.200,0.1434) (3.300,0.1372) (3.400,0.1314) (3.500,0.1259) (3.600,0.1207) (3.700,0.1158) (3.800,0.1112) (3.900,0.1069) (4.000,0.1027)};
\draw[cExact,densely dashed,thick] plot coordinates {(-4.000,0.1000) (-3.900,0.1041) (-3.800,0.1085) (-3.700,0.1131) (-3.600,0.1179) (-3.500,0.1231) (-3.400,0.1285) (-3.300,0.1343) (-3.200,0.1404) (-3.100,0.1470) (-3.000,0.1538) (-2.900,0.1612) (-2.800,0.1689) (-2.700,0.1771) (-2.600,0.1859) (-2.500,0.1951) (-2.400,0.2049) (-2.300,0.2153) (-2.200,0.2262) (-2.100,0.2378) (-2.000,0.2500) (-1.900,0.2628) (-1.800,0.2762) (-1.700,0.2903) (-1.600,0.3049) (-1.500,0.3200) (-1.400,0.3356) (-1.300,0.3515) (-1.200,0.3676) (-1.100,0.3839) (-1.000,0.4000) (-0.900,0.4158) (-0.800,0.4310) (-0.700,0.4454) (-0.600,0.4587) (-0.500,0.4706) (-0.400,0.4808) (-0.300,0.4890) (-0.200,0.4950) (-0.100,0.4988) (0.000,0.5000) (0.100,0.4988) (0.200,0.4950) (0.300,0.4890) (0.400,0.4808) (0.500,0.4706) (0.600,0.4587) (0.700,0.4454) (0.800,0.4310) (0.900,0.4158) (1.000,0.4000) (1.100,0.3839) (1.200,0.3676) (1.300,0.3515) (1.400,0.3356) (1.500,0.3200) (1.600,0.3049) (1.700,0.2903) (1.800,0.2762) (1.900,0.2628) (2.000,0.2500) (2.100,0.2378) (2.200,0.2262) (2.300,0.2153) (2.400,0.2049) (2.500,0.1951) (2.600,0.1859) (2.700,0.1771) (2.800,0.1689) (2.900,0.1612) (3.000,0.1538) (3.100,0.1470) (3.200,0.1404) (3.300,0.1343) (3.400,0.1285) (3.500,0.1231) (3.600,0.1179) (3.700,0.1131) (3.800,0.1085) (3.900,0.1041) (4.000,0.1000)};
\node[cExact,anchor=west] at (1.6,0.36) {極限 $\frac{2}{4+x^2}$};
\end{tikzpicture}}

\newcommand{\EightTabCollapse}{%
\begin{tabular}{@{}lll*{4}{c}@{}}
\toprule
$z_c-z$ & $z$ & $g(z;0)$ & $x=0$ & $x=1$ & $x=2$ & $x=4$\\
\midrule
0.1 & 0.4 & 5 & 0.5000 & 0.4172 & 0.2819 & 0.1317\\
0.05 & 0.45 & 10 & 0.5000 & 0.4085 & 0.2652 & 0.1145\\
0.01 & 0.49 & 50 & 0.5000 & 0.4017 & 0.2529 & 0.1027\\
\midrule
$\to0$ & & & 0.5000 & 0.4000 & 0.2500 & 0.1000\\
\bottomrule
\end{tabular}}

\newcommand{\EightTabXi}{%
\begin{tabular}{@{}cccc@{}}
\toprule
$z/z_c$ & $\xi$（(8.8c)） & $\xi$（格子） & $G_{\text{格子}}/G_{\rm Yukawa}$（$r=8$）\\
\midrule
0.96 & 2.0412 & 2.0205 & 1.0231\\
0.99 & 4.0825 & 4.0722 & 1.0178\\
\bottomrule
\end{tabular}}

\newcommand{\EightFigSD}{%
\begin{tikzpicture}[x=0.16cm,y=2.6cm,>=Stealth,font=\scriptsize,baseline=(current bounding box.center)]
\draw[->] (-21.5,0) -- (22.5,0) node[right] {$x$};
\draw[->] (-21.5,0) -- (-21.5,1.15);
\foreach \x in {-20,-10,10,20} {\draw (\x,0) -- (\x,-0.03); \node[below] at (\x,-0.02) {$\x$};}
\foreach \y in {0.5,1} {\draw (-21.5,\y) -- (-21.9,\y); \node[left] at (-21.9,\y) {$\y$};}
\fill[cBar!35] (-20.8,0) rectangle (-19.2,0.0000);
\fill[cBar!35] (-18.8,0) rectangle (-17.2,0.0001);
\fill[cBar!35] (-16.8,0) rectangle (-15.2,0.0010);
\fill[cBar!35] (-14.8,0) rectangle (-13.2,0.0062);
\fill[cBar!35] (-12.8,0) rectangle (-11.2,0.0262);
\fill[cBar!35] (-10.8,0) rectangle (-9.2,0.0839);
\fill[cBar!35] (-8.8,0) rectangle (-7.2,0.2098);
\fill[cBar!35] (-6.8,0) rectangle (-5.2,0.4196);
\fill[cBar!35] (-4.8,0) rectangle (-3.2,0.6818);
\fill[cBar!35] (-2.8,0) rectangle (-1.2,0.9091);
\fill[cBar!35] (-0.8,0) rectangle (0.8,1.0000);
\fill[cBar!35] (1.2,0) rectangle (2.8,0.9091);
\fill[cBar!35] (3.2,0) rectangle (4.8,0.6818);
\fill[cBar!35] (5.2,0) rectangle (6.8,0.4196);
\fill[cBar!35] (7.2,0) rectangle (8.8,0.2098);
\fill[cBar!35] (9.2,0) rectangle (10.8,0.0839);
\fill[cBar!35] (11.2,0) rectangle (12.8,0.0262);
\fill[cBar!35] (13.2,0) rectangle (14.8,0.0062);
\fill[cBar!35] (15.2,0) rectangle (16.8,0.0010);
\fill[cBar!35] (17.2,0) rectangle (18.8,0.0001);
\fill[cBar!35] (19.2,0) rectangle (20.8,0.0000);
\draw[cExact,very thick] plot coordinates {(-20.00,0.0001) (-19.67,0.0001) (-19.33,0.0001) (-19.00,0.0002) (-18.67,0.0002) (-18.33,0.0003) (-18.00,0.0004) (-17.67,0.0006) (-17.33,0.0008) (-17.00,0.0010) (-16.67,0.0013) (-16.33,0.0017) (-16.00,0.0023) (-15.67,0.0029) (-15.33,0.0037) (-15.00,0.0047) (-14.67,0.0060) (-14.33,0.0075) (-14.00,0.0094) (-13.67,0.0117) (-13.33,0.0145) (-13.00,0.0179) (-12.67,0.0219) (-12.33,0.0267) (-12.00,0.0324) (-11.67,0.0391) (-11.33,0.0470) (-11.00,0.0561) (-10.67,0.0666) (-10.33,0.0787) (-10.00,0.0925) (-9.67,0.1081) (-9.33,0.1257) (-9.00,0.1454) (-8.67,0.1672) (-8.33,0.1914) (-8.00,0.2179) (-7.67,0.2467) (-7.33,0.2779) (-7.00,0.3114) (-6.67,0.3471) (-6.33,0.3848) (-6.00,0.4244) (-5.67,0.4655) (-5.33,0.5080) (-5.00,0.5514) (-4.67,0.5954) (-4.33,0.6395) (-4.00,0.6832) (-3.67,0.7261) (-3.33,0.7676) (-3.00,0.8071) (-2.67,0.8442) (-2.33,0.8784) (-2.00,0.9092) (-1.67,0.9360) (-1.33,0.9586) (-1.00,0.9765) (-0.67,0.9895) (-0.33,0.9974) (0.00,1.0000) (0.33,0.9974) (0.67,0.9895) (1.00,0.9765) (1.33,0.9586) (1.67,0.9360) (2.00,0.9092) (2.33,0.8784) (2.67,0.8442) (3.00,0.8071) (3.33,0.7676) (3.67,0.7261) (4.00,0.6832) (4.33,0.6395) (4.67,0.5954) (5.00,0.5514) (5.33,0.5080) (5.67,0.4655) (6.00,0.4244) (6.33,0.3848) (6.67,0.3471) (7.00,0.3114) (7.33,0.2779) (7.67,0.2467) (8.00,0.2179) (8.33,0.1914) (8.67,0.1672) (9.00,0.1454) (9.33,0.1257) (9.67,0.1081) (10.00,0.0925) (10.33,0.0787) (10.67,0.0666) (11.00,0.0561) (11.33,0.0470) (11.67,0.0391) (12.00,0.0324) (12.33,0.0267) (12.67,0.0219) (13.00,0.0179) (13.33,0.0145) (13.67,0.0117) (14.00,0.0094) (14.33,0.0075) (14.67,0.0060) (15.00,0.0047) (15.33,0.0037) (15.67,0.0029) (16.00,0.0023) (16.33,0.0017) (16.67,0.0013) (17.00,0.0010) (17.33,0.0008) (17.67,0.0006) (18.00,0.0004) (18.33,0.0003) (18.67,0.0002) (19.00,0.0002) (19.33,0.0001) (19.67,0.0001) (20.00,0.0001)};
\draw[cWrong,densely dashed,thick] plot coordinates {(-20.00,0.0000) (-19.67,0.0000) (-19.33,0.0000) (-19.00,0.0000) (-18.67,0.0000) (-18.33,0.0000) (-18.00,0.0000) (-17.67,0.0000) (-17.33,0.0000) (-17.00,0.0000) (-16.67,0.0000) (-16.33,0.0000) (-16.00,0.0000) (-15.67,0.0000) (-15.33,0.0000) (-15.00,0.0000) (-14.67,0.0000) (-14.33,0.0001) (-14.00,0.0001) (-13.67,0.0001) (-13.33,0.0002) (-13.00,0.0003) (-12.67,0.0005) (-12.33,0.0007) (-12.00,0.0011) (-11.67,0.0015) (-11.33,0.0022) (-11.00,0.0031) (-10.67,0.0044) (-10.33,0.0062) (-10.00,0.0085) (-9.67,0.0117) (-9.33,0.0158) (-9.00,0.0211) (-8.67,0.0280) (-8.33,0.0366) (-8.00,0.0475) (-7.67,0.0609) (-7.33,0.0772) (-7.00,0.0970) (-6.67,0.1205) (-6.33,0.1481) (-6.00,0.1801) (-5.67,0.2167) (-5.33,0.2581) (-5.00,0.3041) (-4.67,0.3545) (-4.33,0.4089) (-4.00,0.4668) (-3.67,0.5272) (-3.33,0.5891) (-3.00,0.6514) (-2.67,0.7128) (-2.33,0.7716) (-2.00,0.8266) (-1.67,0.8761) (-1.33,0.9188) (-1.00,0.9535) (-0.67,0.9791) (-0.33,0.9947) (0.00,1.0000) (0.33,0.9947) (0.67,0.9791) (1.00,0.9535) (1.33,0.9188) (1.67,0.8761) (2.00,0.8266) (2.33,0.7716) (2.67,0.7128) (3.00,0.6514) (3.33,0.5891) (3.67,0.5272) (4.00,0.4668) (4.33,0.4089) (4.67,0.3545) (5.00,0.3041) (5.33,0.2581) (5.67,0.2167) (6.00,0.1801) (6.33,0.1481) (6.67,0.1205) (7.00,0.0970) (7.33,0.0772) (7.67,0.0609) (8.00,0.0475) (8.33,0.0366) (8.67,0.0280) (9.00,0.0211) (9.33,0.0158) (9.67,0.0117) (10.00,0.0085) (10.33,0.0062) (10.67,0.0044) (11.00,0.0031) (11.33,0.0022) (11.67,0.0015) (12.00,0.0011) (12.33,0.0007) (12.67,0.0005) (13.00,0.0003) (13.33,0.0002) (13.67,0.0001) (14.00,0.0001) (14.33,0.0001) (14.67,0.0000) (15.00,0.0000) (15.33,0.0000) (15.67,0.0000) (16.00,0.0000) (16.33,0.0000) (16.67,0.0000) (17.00,0.0000) (17.33,0.0000) (17.67,0.0000) (18.00,0.0000) (18.33,0.0000) (18.67,0.0000) (19.00,0.0000) (19.33,0.0000) (19.67,0.0000) (20.00,0.0000)};
\node[cExact,anchor=west] at (8.5,0.85) {更正 $e^{-x^2/42}$};
\node[cWrong,anchor=west] at (8.5,0.7) {書上 $e^{-x^2/21}$};
\node[cBar,anchor=west] at (8.5,0.55) {精確（長條）};
\end{tikzpicture}}

\newcommand{\EightTabSD}{%
\begin{tabular}{@{}rrccc@{}}
\toprule
$N$ & $x$ & 精確 & 更正 $e^{-x^2/2(N+1)}$ & 書上 $e^{-x^2/(N+1)}$\\
\midrule
20 & 4 & 0.6818 & 0.6832 & 0.4668\\
20 & 8 & 0.2098 & 0.2179 & 0.0475\\
100 & 8 & 0.7282 & 0.7285 & 0.5306\\
100 & 12 & 0.4894 & 0.4902 & 0.2403\\
\bottomrule
\end{tabular}}

\newcommand{\EightTabFisher}{%
\begin{tabular}{@{}lcccc@{}}
\toprule
$\nu$ & $\frac{1}{2}$ & $\frac{3}{5}$ & $\frac{3}{4}$ & 0.588\\
$\frac1{1-\nu}$ & $2$ & $\frac{5}{2}$ & $4$ & 2.43\\
\midrule
 & 高斯 & 三維 Flory & 二維自迴避 & 三維自迴避\\
\bottomrule
\end{tabular}}

\newcommand{\EightTabRatio}{%
\begin{tabular}{@{}l*{4}{c}c@{}}
\toprule
$N$ & 10 & 12 & 14 & 16 & 已知值\\
\midrule
$\hat\mu_N$ & 2.6592 & 2.6529 & 2.6483 & 2.6454 & 2.638\\
$\hat\gamma_N$ & 1.3552 & 1.3519 & 1.3496 & 1.3483 & $\frac{43}{32}$\\
\bottomrule
\end{tabular}}

\newcommand{\EightTabAmp}{%
\begin{tabular}{@{}l*{5}{c}@{}}
\toprule
$d$ & 1 & 2 & 3 & 4 & 5\\
\midrule
$\frac{4-d}{2}$ & $\frac{3}{2}$ & $1$ & $\frac{1}{2}$ & $0$ & $-\frac{1}{2}$\\
$N^{(4-d)/2}$（$N=100$） & 1000 & 100 & 10 & 1 & 0.1\\
\bottomrule
\end{tabular}}

\newcommand{\EightTabNu}{%
\begin{tabular}{@{}cccl@{}}
\toprule
$d$ & 放大鏡 $\frac3{d+2}$ & 第 7 章拔河 $\nu_F$ & 已知值\\
\midrule
1 & $1$ & $1$ & 1\\
2 & $\frac{3}{4}$ & $\frac{3}{4}$ & $\frac34$（精確）\\
3 & $\frac{3}{5}$ & $\frac{3}{5}$ & 0.588（數值）\\
4 & $\frac{1}{2}$ & $\frac{1}{2}$ & $\frac12$\\
\bottomrule
\end{tabular}}

\newcommand{\EightTabErrata}{%
\begin{tabular}{@{}lll@{}}
\toprule
位置 & 書上 & 應為\\
\midrule
p.~195，(8.6)$\to$(8.7) & {\color{cWrong}replace $k$ by $w=k/r$} & {\color{cExact}$w=kr$（$k=w/r$）}\\
p.~196，第一段 & {\color{cWrong}$z_c$ is the coordination number} & {\color{cExact}$z_c=1/(\text{配位數})=\frac1{2d}$}\\
p.~198，(8.18) & {\color{cWrong}$\nu|\vec r|(z_c-z)^{\nu-1}A'-\frac{N+1}z=0$} & {\color{cExact}$-\nu|\vec r|(z_c-z)^{\nu-1}A'-\frac{N+1}z=0$}\\
p.~198，(8.20) 第一行 & {\color{cWrong}$\exp\big(-B|\vec r|\delta^\nu-\frac{N+1}{z_c}\big)$} & {\color{cExact}$z_c^{-(N+1)}\exp\big(-B|\vec r|\delta^\nu+\frac{(N+1)\delta}{z_c}\big)$}\\
p.~198，(8.20) 第二行 & {\color{cWrong}$\exp[-((\nu z_c)^\nu B)^{\frac1{1-\nu}}\cdots]$} & {\color{cExact}$\exp[-(1-\nu)((\nu z_c)^\nu B)^{\frac1{1-\nu}}\cdots]$}\\
p.~200，(8.22) 前 & {\color{cWrong}$\vec r\,'(s')=l_1\vec r(l_2s)$} & {\color{cExact}$\vec r\,'(s')=l_1\vec r(s'/l_2)$}\\
p.~200，Ex.~8.6 & {\color{cWrong}rescaled by $l_1^d/l_2$} & {\color{cExact}$l_1^d/l_2^2$}\\
\bottomrule
\end{tabular}}


\section{標度入門（第 8 章）}

\begin{keybox}[先看這裡]
前七章的結論幾乎都是高斯：走 $N$ 步，寬度 $\propto\sqrt N$。這一章問兩件事：\textbf{什麼時候不是高斯？}以及\textbf{不是高斯時，還剩下什麼規律？}
第一個答案：步長的分布尾巴太肥（二階動差發散）時，中央極限定理不管用，寬度長得比 $\sqrt N$ 快、分布的尾巴是冪次。第二個答案是\textbf{標度}：
靠近臨界逸度 $z_c$ 時只剩一個長度 $\xi\propto(z_c-z)^{-\nu}$，用它當尺，所有曲線疊成一條；由標度形式用最陡下降法可以讀出
$C(N;\vec r)\sim\exp[-(r/N^\nu)^{1/(1-\nu)}]$——只有 $\nu=\frac12$ 才是高斯。最後用「放大鏡」（重新標度）重新推出第 7 章 Flory 的 $\nu=\frac3{d+2}$。本章只有 10 頁，卻是第 9--12 章（臨界現象、重整化群）的入口。

\textbf{記號}：步數 $N$、條數 $C(N;\vec r)$、一步的結構函數 $\chi(\vec q)$、生成函數 $g(z;\vec q)$（傅立葉）與 $G(z;\vec r)$（實空間）、逸度 $z$（簡單立方 $z_c=\frac1{2d}$）、
標度指數 $\nu,\eta$（普通漫步 $\frac12$、0）、$\gamma$、連接常數 $\mu$。長尾巴步伐：步長 $r\ge1$、$P(\text{步長}>r)=r^{-y}$（三維空間密度 $\propto r^{-3-y}$，書上寫 $r^{-5+\delta}$，即 $y=2-\delta$）。
8.2 節的 $a$ 是 $\chi(\vec 0)$（簡單立方 $a=2d$、$b=1$），\textbf{不是格距}。重新標度：$s'=l_2s$、$\vec r\,'=l_1\vec r$。

\textbf{全章的主例}：長尾巴（$\delta=1$）1000 步的寬度 1822，步長 1 只有 28.06；一維塌縮 $(z_c-z)g\to\frac2{4+x^2}$；一維 $N=20$、$x=8$ 的比值：精確 0.2098、更正的 (8.20) 0.2179、書上 0.0475。
\begin{center}
\small
\begin{tabular}{@{}lll@{\qquad}lll@{}}
\toprule
節 & 內容 & 動畫 & 節 & 內容 & 動畫\\
\midrule
8.1--8.1.1 & 普適性的例外、長尾巴 & W801 & 8.1.2（後半） & 最陡下降、$\gamma$、比值法 & W803\\
8.1.2（前半） & 生成函數的標度形式 & W802 & 8.1.3 & 重新標度、Flory、回顧 & W804\\
\bottomrule
\end{tabular}
\end{center}
（上表是書上的節號；本講義 8.1--8.4 依序對應。）數字一律由 \texttt{codes/utils/walk\_scale.py} 重算；逐式校對見 \texttt{ref/md/ch08\_formulas.md}；
圖表由 \texttt{ch8/make\_figs.py} 產生並先做數字檢查。連續播放：\texttt{Ch8Demo}（W801--W804）。
\end{keybox}

% =====================================================================
\subsection{普適性與它的例外：長尾巴的步伐}
\leavevmode\bk{pp.~193--195}\ \anim{W801}\par\nopagebreak

\textbf{問}：如果醉漢偶爾來一個超大步，走 1000 步之後，還是一團高斯的霧嗎？

\textbf{直觀}：第 2 章的高斯只用到一件事——一步的結構函數能展開成 $\chi(\vec q)=A-B|\vec q|^2+\cdots$，而 $B$ 就是步長的二階動差。
步長分布的細節（格子、連續、每步長短不一）全部被忘掉，只剩 $B$：這就是\textbf{普適性}。反過來，二階動差一發散，$\chi$ 就不能展開到 $q^2$，高斯也就沒了。

\textbf{二階動差發散}（(8.1)(8.2)）：三維裡步長的密度 $p(|\vec r|)\propto|\vec r|^{-\alpha}$，則
\[\int p(|\vec r|)\,|\vec r|^2\,d\vec r\ \propto\ \int|\vec r|^{4-\alpha}\,d|\vec r|,\tag{8.2}\]
$\alpha\le5$ 時在大 $r$ 發散。〔注〕書上寫「$\alpha$ 小於 5」；$\alpha=5$ 是對數發散，也不行。
\textbf{算到底}：本講義（與動畫）的主例 $P(\text{步長}>r)=r^{-1}$（$y=1$、$\alpha=4$、$\delta=1$），徑向密度 $r^{-2}$；截到 $R$ 的二階動差
$\langle r^2\rangle_R=\int_1^Rr^{-2}\,r^2\,dr=R-1$：$R=10,100,1000$ 得 9、99、999——截得越遠越大，沒有極限。

\textbf{結構函數不能展開}（(8.3)(8.4)）：$p\propto|\vec r|^{-5+\delta}$（$0<\delta<2$）時
\[\chi(\vec q)\approx A-B|\vec q|^{2-\delta}.\tag{8.4}\]
主例 $\delta=1$：$\chi(q)=1-\frac\pi4q+\cdots$，$B=\frac\pi4=0.7854$；數值積分 $\frac{1-\chi(q)}q=0.7687,\ 0.7837,\ 0.7852$（$q=0.1,0.01,0.001$），
$\chi(0.01)=0.99216$ 對 $1-\frac\pi4(0.01)=0.99215$。在對數圖上 $1-\chi$ 的斜率是 $2-\delta$：$\delta=1$ 量到 1.00、$\delta=\frac12$ 量到 1.49，步長 1 的步伐是 2（圖 \ref{eight-fig-levy} 左）。

\textbf{寬度不是 $\sqrt N$}：三維、各 2000 人（種子 801），$|R_N|$ 的中位數（圖 \ref{eight-fig-levy} 右上）。長尾巴 $N=10,100,1000$：18、172、1822，除以 $N$ 是 1.77、1.72、1.82，
趨近理論 $\frac\pi4\times2.2644=1.7785$（$\chi^N\approx e^{-NBq}$ 正是三維 Cauchy 分布，2.2644 是它的徑向中位數）；步長 1：2.86、8.98、28.06，除以 $\sqrt N$ 是 0.90、0.90、0.89，
趨近 $1.5382/\sqrt3=0.8881$。log--log 斜率 1.01 對 0.50：寬度 $\propto N^{1/(2-\delta)}$，主例長得像 $N$。

\textbf{冪次的尾巴}（(8.5)）：
\[C(N;\vec x-\vec y)=\int e^{i\vec q\cdot(\vec x-\vec y)}A^Ne^{-N(B/A)|\vec q|^{2-\delta}}d\vec q\ \propto\ |\vec x-\vec y|^{-5+\delta}.\tag{8.5}\]
數值逆 Fourier 除以換變數得的尾巴（$NB/A=1$）：$\delta=1$ 在 $r=10,30,100$ 是 0.9803、0.9978、0.9998；$\delta=\frac12$ 是 1.1755、1.0316、1.0051；
$r=100$ 的局部斜率 $-4.00$、$-4.51$，就是 $-5+\delta=-4,\ -\frac92$（圖 \ref{eight-fig-levy} 右下）。$\delta=1$ 就是第 2 章 Cauchy 的尾巴 $\frac{N}{\pi^2r^4}$（(2.31)）。
〔注〕(8.5) 說的是尾巴（$r\gg N^{1/(2-\delta)}$），前面還有一個因子 $N$（$N$ 加倍，尾巴加倍：遠處多半是「一大跳」造成的）；一般 $d$ 維是 $r^{-(d+2-\delta)}$。

\textbf{換變數的技巧}（(8.6)(8.7)）：想知道 $\int e^{ikr}k^y\,dk$ 怎麼隨 $r$ 變，令 $w=kr$（$k=w/r$、$dk=dw/r$）：
\[\int e^{ikr}k^y\,dk=r^{-(y+1)}\int e^{iw}w^y\,dw .\tag{8.6, 8.7}\]
右邊的積分是一個與 $r$ 無關的數，$r$ 的冪次一眼讀出。\textbf{核對}：$y=\frac12$，$\big|\int_0^\infty e^{iw}w^{1/2}dw\big|=\Gamma(\frac32)=0.8862$；
數值上 $r=2$ 與 $r=1$ 的積分之比 0.3536，正是 $2^{-3/2}$。練習 8.1 就是用這一招得到 $r^{-5+\delta}$。
\begin{errbox}[p.~195，(8.6)$\to$(8.7)]
書上說「replace $k$ by $w=k/r$」；應為 $w=kr$。照書上 $k=rw$，積分變成 $r^{y+1}\int e^{ir^2w}w^y\,dw$，$r$ 還留在指數裡，提不出來。
\end{errbox}

\begin{figure}[htbp]
\centering
\begin{minipage}[c]{0.46\linewidth}\centering\EightFigChi\end{minipage}\hfill
\begin{minipage}[c]{0.52\linewidth}\centering\footnotesize\setlength{\tabcolsep}{4pt}\EightTabWidth\\[6pt]\EightTabTail\end{minipage}
\caption{長尾巴的步伐（三維）。左：$1-\chi(q)$ 的對數圖，斜率 $2-\delta$（點是數值積分；虛線是步長 1 的 $1-\frac{\sin q}q$，斜率 2）。
右上：$|R_N|$ 的中位數，長尾巴 $\propto N$、步長 1 $\propto\sqrt N$。右下：(8.5) 數值逆 Fourier 與冪次尾巴之比，及 $r=100$ 的局部斜率}\label{eight-fig-levy}
\end{figure}

\textbf{揭曉}：會跳的醉漢不是一團高斯的霧，而是幾團霧被長跳連起來；寬度 $\propto N^{1/(2-\delta)}$，尾巴是冪次 $r^{-5+\delta}$。

% =====================================================================
\subsection{生成函數的標度形式}
\leavevmode\bk{pp.~195--197}\ \anim{W802}\par\nopagebreak

\textbf{問}：$z$ 越靠近 $z_c$，$g(z;q)$ 越尖越高；不同 $z$ 的曲線，能不能換個座標疊成一條？

\textbf{靠近 $z_c$ 的生成函數}（(8.8)）：普通的步伐 $\chi(\vec q)=a-b|\vec q|^2+O(|\vec q|^4)$，所以
\[g(z;\vec q)=\frac1{1-z\chi(\vec q)}\approx\frac1{1-az+zb|\vec q|^2}\approx\frac1a\,\frac1{z_c-z+(b/a^2)|\vec q|^2},\qquad z_c=\frac1a=\frac1{2d}.\tag{8.8}\]
\begin{errbox}[p.~196，第一段]
書上說 $z_c$「turns out to be the coordination number」；應為配位數的\textbf{倒數}，$z_c=\frac1{2d}$（同頁 (8.8c) 後面也寫 $1/2d$）。
照書上一維 $z_c=2$，可是 $z>\frac12$ 時 $\sum_N(2z)^N$ 早就發散了。
\end{errbox}

\textbf{算到底}（一維，$a=2$、$b=1$，$g=\frac1{1-2z\cos q}$）：$z=0.4,\ 0.45,\ 0.49$（$z_c-z=0.1,\ 0.05,\ 0.01$）時 $g(z;0)=\frac1{1-2z}=5,\ 10,\ 50$，曲線越來越尖（圖 \ref{eight-fig-collapse} 左）。
記 $\varepsilon=z_c-z$，令 $q=x\sqrt\varepsilon$：$1-2z\cos q\approx(1-2z)+zq^2\approx2\varepsilon+\frac12\varepsilon x^2$，所以
\[(z_c-z)\,g\ \approx\ \frac{\varepsilon}{2\varepsilon+\frac12\varepsilon x^2}=\frac2{4+x^2},\]
$\varepsilon$ 消掉了。以 $x=2$ 為例：三個 $z$ 的 $(z_c-z)g$ 是 0.2819、0.2652、0.2529，趨近 $\frac2{4+4}=0.25$。三條曲線疊成一條（圖 \ref{eight-fig-collapse} 中）——這叫\textbf{資料塌縮}。

\textbf{標度形式}（(8.9)(8.10)）：把 (8.8) 改寫成
\[g(z;\vec q)=|\vec q|^{-2}f\big(q(z_c-z)^{-1/2}\big),\tag{8.9}\]
意思是：$g$ 只透過「$q$ 和 $(z_c-z)^{1/2}$ 的比」依賴 $z$。普遍的假設是把兩個 $\frac12$、$-2$ 換成未知的指數：
\[g(z;\vec q)=|\vec q|^{-2+\eta}f\big(q(z_c-z)^{-\nu}\big)\qquad(\text{普通漫步 }\nu=\tfrac12,\ \eta=0;\ \text{自迴避漫步不同}).\tag{8.10}\]

\textbf{回到實空間}（(8.11)--(8.14)）：$G(z;\vec r)=\int g(z;\vec q)e^{i\vec q\cdot\vec r}d^dq$，令 $\vec q=\vec\kappa/|\vec r|$，$d^dq$ 給 $|\vec r|^{-d}$、$|\vec q|^{-2+\eta}$ 給 $|\vec r|^{2-\eta}$：
\[G(z;\vec r)=|\vec r|^{-d+2-\eta}F\big(|\vec r|(z_c-z)^\nu\big),\qquad F(x)\equiv\int|\vec\kappa|^{-2+\eta}f(\kappa/|x|)\,e^{i\vec\kappa\cdot\vec x/|\vec x|}d^d\kappa .\tag{8.13, 8.14}\]
〔注〕(8.11) 省略了 $(2\pi)^{-d}$。書上假設 $F$ 在 $x\to0$ 趨於常數、$x\to\infty$ 指數衰減；前者只在 $d>2$ 成立（一維 $F(u)=e^{-2u}/2u$，$u\to0$ 發散）。

\textbf{三維實例}（簡單立方 $a=6$、$b=1$）：(8.8) 的逆 Fourier 是 Yukawa 形式
\[G(z;r)\approx\frac3{2\pi}\,\frac{e^{-r/\xi}}r,\qquad\xi=\frac1{\sqrt{6(1-z/z_c)}}\propto(z_c-z)^{-1/2},\]
即 $rG=F(u)=\frac3{2\pi}e^{-6u}$、$u=r\sqrt{z_c-z}$，$F(0)=\frac3{2\pi}=0.4775$——兩個假設都滿足。$z/z_c=0.96$：$\xi=2.0412$（格子上的精確衰減長度 2.0205）；
$z/z_c=0.99$：$\xi=4.0825$（4.0722）。$z_c-z$ 變成 $\frac14$，$\xi$ 變成 2 倍；$r=8$ 處格子精確 $G$ 與 Yukawa 之比 1.0231、1.0178（圖 \ref{eight-fig-collapse} 下）。

\begin{figure}[htbp]
\centering
\EightFigRawG\hfill\EightFigCollapse\\[6pt]
{\footnotesize\setlength{\tabcolsep}{3.5pt}\EightTabCollapse\hfill\EightTabXi}
\caption{標度形式。左上：一維 $g(z;q)=\frac1{1-2z\cos q}$，$z=0.4,0.45,0.49$。右上：縱軸乘 $z_c-z$、橫軸用 $x=q/\sqrt{z_c-z}$，三條曲線疊在虛線 $\frac2{4+x^2}$ 上。
左下：塌縮的數值。右下：三維的關聯長度 $\xi$，(8.8c) 與格子精確值}\label{eight-fig-collapse}
\end{figure}

\textbf{揭曉}：靠近 $z_c$ 只剩一個長度 $\xi\propto(z_c-z)^{-\nu}$；用 $\xi$ 當尺，所有曲線疊成一條。

% =====================================================================
\subsection{最陡下降：從標度形式讀出分布}
\leavevmode\bk{pp.~197--199}\ \anim{W803}\par\nopagebreak

\textbf{問}：知道 $G(z;\vec r)$ 的標度形式，怎麼拿回「$N$ 步走到 $\vec r$」的條數？

\textbf{Cauchy 公式＋鞍點}（(8.15)--(8.17)）：$C(N;\vec r)$ 是 $G$ 裡 $z^N$ 的係數，
\[C(N;\vec r)=\frac1{2\pi i}\oint\frac{G(z;\vec r)}{z^{N+1}}dz\approx|\vec r|^{-d+2-\eta}F\big(|\vec r|(z_c-z)^\nu\big)z^{-(N+1)}\Big|_{z\text{ 取極值}} .\tag{8.15, 8.16}\]
要極值化的是 $A\big(|\vec r|(z_c-z)^\nu\big)-(N+1)\ln z$，$A=\ln F$；$F$ 指數衰減，所以 $x\gg1$ 時 $A(x)\approx-Bx$。對 $z$ 微分（$\frac{d}{dz}(z_c-z)^\nu=-\nu(z_c-z)^{\nu-1}$）：
\[-\nu|\vec r|(z_c-z)^{\nu-1}A'-\frac{N+1}z=0\quad\Rightarrow\quad\delta\equiv z_c-z\approx\Big(\frac{N+1}{\nu Bz_c|\vec r|}\Big)^{1/(\nu-1)}.\tag{8.18, 8.19}\]
\begin{errbox}[p.~198，(8.18)]
第一項沒有負號：$\nu|\vec r|(z_c-z)^{\nu-1}A'-\frac{N+1}z=0$。照書上，$A'=-B$ 時兩項都是負的，沒有解。
\end{errbox}
代回去：$-(N+1)\ln(z_c-\delta)\approx-(N+1)\ln z_c+\frac{(N+1)\delta}{z_c}$，而 (8.19) 說 $B|\vec r|\delta^\nu=\frac{(N+1)\delta}{\nu z_c}$，兩項合起來是 $-\frac{(1-\nu)(N+1)\delta}{\nu z_c}$：
\[C(N;\vec r)\approx z_c^{-(N+1)}e^{-B|\vec r|\delta^\nu+(N+1)\delta/z_c}\propto\exp\Big[-(1-\nu)\big((\nu z_c)^\nu B\big)^{\frac1{1-\nu}}(N+1)^{-\frac{\nu}{1-\nu}}|\vec r|^{\frac1{1-\nu}}\Big],\tag{8.20}\]
也就是 $C\sim\exp\big[-c\,(r/N^\nu)^{1/(1-\nu)}\big]$。\textbf{檢驗}：簡單立方 $\nu=\frac12$、$B=2d$、$z_c=\frac1{2d}$，係數 $\frac12\cdot\big(\sqrt{1/4d}\cdot2d\big)^2=\frac d2$，
指數 $-\frac{d\,r^2}{2N}$——正是第 2 章的高斯 (2.20)（$d=3$：係數 $\frac32$）。
\begin{errbox}[p.~198，(8.20)]
第一行印成 $\exp\big(-B|\vec r|\delta^\nu-(N+1)/z_c\big)$，應為 $z_c^{-(N+1)}\exp\big(-B|\vec r|\delta^\nu+(N+1)\delta/z_c\big)$；
第二行的指數少了因子 $(1-\nu)$。照書上 $\nu=\frac12$ 得 $e^{-dr^2/N}$，變異數少一半（一維 $e^{-x^2/N}$）。
\end{errbox}

\textbf{算到底}（一維 $N=20$、$x=8$；$\nu=\frac12$、$B=2$、$z_c=\frac12$）：(8.19) 給 $\delta=\big(\frac{21}{\frac12\cdot2\cdot\frac12\cdot8}\big)^{-2}=\big(\frac{21}4\big)^{-2}=\frac{16}{441}$，$\delta^{1/2}=\frac4{21}$。
指數的兩項：$-B x\delta^{1/2}=-2\cdot8\cdot\frac4{21}=-\frac{64}{21}$，$+\frac{(N+1)\delta}{z_c}=21\cdot\frac{16}{441}\cdot2=+\frac{32}{21}$，和是 $-\frac{32}{21}=-\frac{64}{42}$，
即 $e^{-x^2/2(N+1)}=e^{-64/42}=0.2179$；精確的 $\binom{20}{14}/\binom{20}{10}=0.2098$。照書上只剩第一項：$e^{-64/21}=0.0475$，差了四倍多（圖 \ref{eight-fig-sd}）。
$N=100$ 時更正版和精確值幾乎一樣（$x=8$：0.7285 對 0.7282；$x=12$：0.4902 對 0.4894）。〔注〕一維 $G$ 在 $z=\pm\frac12$ 各有一個鞍點，完整的條數要乘 2（$N+x$ 偶數時），取比值時消掉。

\begin{figure}[htbp]
\centering
\begin{minipage}[c]{0.5\linewidth}\centering\EightFigSD\end{minipage}\hfill
\begin{minipage}[c]{0.48\linewidth}\centering\footnotesize\setlength{\tabcolsep}{3pt}\EightTabSD\end{minipage}\\[8pt]
{\footnotesize\setlength{\tabcolsep}{4pt}\EightTabFisher\hfill\EightTabRatio}
\caption{最陡下降。左上：一維 $N=20$ 的 $C(20;x)/C(20;0)$（長條），更正的 (8.20)（綠）與書上的 (8.20)（紅虛線）。右上：$N=20$、100 的比較。
左下：指數 $\frac1{1-\nu}$（Fisher）。右下：二維自迴避 $c_N$ 的比值法（$c_{15},c_{16}$ 是文獻值）}\label{eight-fig-sd}
\end{figure}

\textbf{不是高斯}：$C\sim e^{-(r/N^\nu)^{1/(1-\nu)}}$ 的指數 $\frac1{1-\nu}$（Fisher）：$\nu=\frac12\to2$（高斯）、$\frac35\to\frac52$、二維自迴避 $\frac34\to4$；三維 $\nu\approx0.588$（數值，只引用）$\to2.43$。
自迴避漫步的分布在遠處掉得比高斯快得多。

\textbf{總條數與 $\gamma$}（練習 8.3、8.4）：對 $\vec r$ 求和就是 $g(z;\vec 0)$，(8.10) 給 $g(z;0)\propto(z_c-z)^{-\gamma}$，$\gamma=(2-\eta)\nu$。普通漫步 $(2-0)\cdot\frac12=1$；
二維自迴避 $\eta=\frac5{24}$、$\nu=\frac34$ 得 $\gamma=\frac{43}{32}$（已知值，只引用）。取 $z^N$ 的係數：條數 $\propto x^NN^p$，$x=\frac1{z_c}$、$p=\gamma-1$；
二維自迴避 $x=\mu\approx2.638$、$p=\frac{11}{32}=0.34375$。\textbf{和精確計數比}（比值法，第 7 章的 $c_N$）：$\hat\mu_N=\frac{Nr_N-(N-2)r_{N-2}}2$（$r_N=c_N/c_{N-1}$）、
$\hat\gamma_N=1+\frac N2\big(\frac{c_N}{c_{N-2}\mu^2}-1\big)$；$N=10,12,14,16$：$\hat\mu=2.6592,\ 2.6529,\ 2.6483,\ 2.6454\to2.638$，$\hat\gamma=1.3552,\ 1.3519,\ 1.3496,\ 1.3483\to1.34375$。

\textbf{揭曉}：$\nu=\frac12$ 才是高斯；$\nu\ne\frac12$ 時分布是 $\exp[-(r/N^\nu)^{1/(1-\nu)}]$。

% =====================================================================
\subsection{放大鏡：重新標度與 Flory}
\leavevmode\bk{pp.~199--202}\ \anim{W804}\par\nopagebreak

\textbf{問}：把漫步「縮小來看」——步數變少、長度變短——自迴避的效果會變大還是變小？

\textbf{路徑積分}（(8.21)，第 7 章的 (7.16)）：指數裡有兩項，動能項與自迴避項，
\[C(N;\vec x,\vec y)\propto\prod_s\int d\vec r(s)\exp\Big[-\frac3{2a^2}\int_0^N\Big|\frac{d\vec r}{ds}\Big|^2ds-\frac{v_0}2\int_0^N\!\!\int_0^Nds_1ds_2\,\delta\big(\vec r(s_1)-\vec r(s_2)\big)\Big].\tag{8.21}\]
\textbf{放大鏡}：$s'=l_2s$、$\vec r\,'(s')=l_1\vec r(s'/l_2)$（$l_1,l_2<1$：縮小）。動能項：$\frac{d\vec r}{ds}=\frac{l_2}{l_1}\frac{d\vec r\,'}{ds'}$、$ds=ds'/l_2$，乘上 $\frac{l_2}{l_1^2}$；
自迴避項：$ds_1ds_2=ds_1'ds_2'/l_2^2$、$\delta(\vec r\,'/l_1)=l_1^d\delta(\vec r\,')$，乘上 $\frac{l_1^d}{l_2^2}$：
\[-\frac3{2a^2}\frac{l_2}{l_1^2}\int_0^{l_2N}\Big|\frac{d\vec r\,'}{ds'}\Big|^2ds'-\frac{v_0}2\frac{l_1^d}{l_2^2}\int_0^{l_2N}\!\!\int_0^{l_2N}ds_1'ds_2'\,\delta\big(\vec r\,'(s_1')-\vec r\,'(s_2')\big).\tag{8.22}\]
\textbf{核對}：$d=3$、$l_1=\frac12$、$l_2=\frac14$：$\frac{l_2}{l_1^2}=1$、$\frac{l_1^3}{l_2^2}=\frac{1/8}{1/16}=2$；離散高斯鏈的數值量到 1.00、2.00。
\begin{errbox}[p.~200，(8.22) 前]
書上寫 $\vec r(s)\to\vec r\,'(s')=l_1\vec r(l_2s)$；應為 $l_1\vec r(s)=l_1\vec r(s'/l_2)$：新曲線在 $s'$ 的位置，是舊曲線在 $s=s'/l_2$ 的位置放大 $l_1$ 倍。
\end{errbox}

\textbf{上臨界維度 4}（練習 8.6、8.7）：取 $l_1=\sqrt{l_2}$，動能項不變（步長 $a'=a$），自迴避項乘上 $l_2^{(d-4)/2}$。把漫步縮到只剩一步（$l_2=\frac1N$）：
沒有自迴避時那一步長 $a$，放大回去 $R=a/l_1=a\sqrt N$；自迴避卻被放大 $N^{(4-d)/2}$ 倍。$N=100$、$d=1..5$：1000、100、10、1、0.1（圖 \ref{eight-fig-flory} 左）。
$d<4$ 越縮越重要，$d>4$ 越縮越不重要——和第 7 章微擾的冪次表（$d=4$ 那列全是 1）同一個分界。
\begin{errbox}[p.~200，練習 8.6]
說自迴避項「rescaled by a factor $l_1^d/l_2$」；應為 $l_1^d/l_2^2$（同頁正文與 (8.22)）。照書上會得 $l_2^{(d-2)/2}$，上臨界維度變成 2。
〔注〕題目說的「first term in the double integral」指指數裡的第一項（動能項）。
\end{errbox}

\textbf{讓兩項一起縮放}（(8.23)--(8.25)）：要兩項保持同樣的比例，
\[\frac{l_2}{l_1^2}=\frac{l_1^d}{l_2^2}\quad\Rightarrow\quad l_1=l_2^{3/(d+2)};\qquad l_2=\frac1N:\ R\sim l_1^{-1}=N^{3/(d+2)},\quad\nu=\frac3{d+2}.\tag{8.24, 8.25}\]
\textbf{算到底}：$d=3$、$l_2=0.01$（100 步縮成 1 步）：$l_1=0.01^{3/5}=0.0631$，$1/l_1=15.85=100^{3/5}$。$d=1..4$：$\nu=1,\frac34,\frac35,\frac12$，
和第 7 章 Flory 的「拔河」一模一樣（圖 \ref{eight-fig-flory} 右）：一次是熵與排斥拔河，一次是放大鏡底下兩項同比例，兩條路得到同一個 $\frac3{d+2}$。

\begin{figure}[htbp]
\centering
{\footnotesize\setlength{\tabcolsep}{4pt}\EightTabAmp\hfill\EightTabNu}
\caption{重新標度。左：$l_1=\sqrt{l_2}$、$l_2=\frac1N$ 時自迴避項的放大倍數 $N^{(4-d)/2}$，上臨界維度 4。右：兩項同比例得到的 $\nu=\frac3{d+2}$、第 7 章 Flory 的拔河與已知值}\label{eight-fig-flory}
\end{figure}

\textbf{揭曉}：放大鏡底下，四維以下自迴避越看越重要；讓兩項一起縮放，就得到 Flory 的 $\frac3{d+2}$。

% =====================================================================
\subsection{練習 8.1--8.8 解答}
{\small
\begin{enumerate}[label=\textbf{8.\arabic*},itemsep=0.05em]
\item 冪次的尾巴。\quad 答：三維 $\int d^3q\,e^{i\vec q\cdot\vec r}f(q)=\frac{4\pi}r\int_0^\infty q\sin(qr)f(q)\,dq$。遠處把 $A^Ne^{-N(B/A)q^{2-\delta}}\approx A^N\big(1-N\frac BAq^{2-\delta}\big)$：
第一項只給 $\vec r=0$；第二項 $\frac{4\pi}r\int q^{3-\delta}\sin(qr)\,dq$，令 $w=qr$ 得 $r^{-1}\cdot r^{-(4-\delta)}=r^{-5+\delta}$。含 $(2\pi)^{-3}$ 的係數是
$\frac{NB}{A}A^N\frac{\Gamma(4-\delta)\sin(\pi(2-\delta)/2)}{2\pi^2}$；$\delta=1$ 即 Cauchy 的 $\frac{N}{\pi^2r^4}$（數值比 0.9803、0.9978、0.9998）。
\item 最陡下降。\quad 答：$\frac{d}{dz}\big[A(|\vec r|\delta^\nu)-(N+1)\ln z\big]=-\nu|\vec r|\delta^{\nu-1}A'-\frac{N+1}z=0$，$A'=-B$、$z\approx z_c$ 得 (8.19)。代回：$-Br\delta^\nu=-\frac{(N+1)\delta}{\nu z_c}$，
加上 $+\frac{(N+1)\delta}{z_c}$，指數 $-(1-\nu)\big((\nu z_c)^\nu B\big)^{1/(1-\nu)}(N+1)^{-\nu/(1-\nu)}r^{1/(1-\nu)}$。$\nu=\frac12$、簡單立方得 $-\frac{dr^2}{2N}$；
書上少了 $1-\nu$，得 $-\frac{dr^2}N$（一維 $N=20$、$x=8$：0.0475，精確 0.2098）。
\item $\gamma=(2-\eta)\nu$。\quad 答：$\sum_{\vec r}G(z;\vec r)=g(z;\vec 0)$。(8.10) 若一直用到 $q\to0$，而 $z<z_c$ 時 $g(z;0)$ 有限，必須 $f(x)\approx c\,x^{2-\eta}$（$x\to0$），
於是 $g(z;0)=|\vec q|^{-2+\eta}c\,q^{2-\eta}(z_c-z)^{-(2-\eta)\nu}=c\,(z_c-z)^{-(2-\eta)\nu}$，$\gamma=(2-\eta)\nu$（普通 1；二維自迴避 $\frac{43}{32}$）。
\item 總條數。\quad 答：用第 2 章的 (2:S3-10)，$(z_c-z)^{-\gamma}$ 裡 $z^N$ 的係數 $\approx\frac{N^{\gamma-1}}{\Gamma(\gamma)}z_c^{-N-\gamma}$，所以條數 $\propto x^NN^p$，$x=\frac1{z_c}$、$p=\gamma-1$。
普通漫步 $x=2d$、$p=0$；二維自迴避 $x=\mu\approx2.638$、$p=\frac{11}{32}$（$N=1000\to2000$ 的局部斜率 0.3436）；比值法見圖 \ref{eight-fig-sd}。
\item 小 $r$ 的條數。\quad 答：固定 $r$、$N\to\infty$ 時條數應 $\sim z_c^{-N}N^{-d\nu}$（終點散在 $N^{d\nu}$ 的體積裡）。$F(x)=F_0+F_1x^s+\cdots$ 的非解析項給 $(z_c-z)^{\nu s}$，
它的 $z^N$ 係數 $\propto N^{-\nu s-1}$；令 $=N^{-d\nu}$ 得 $s=d-\frac1\nu$。$r$ 的依賴 $r^{-d+2-\eta+s}=r^\theta$，$\theta=\frac{\gamma-1}\nu$。普通漫步 $\theta=0$、
$C(N;0)\propto(2d)^NN^{-d/2}$（$d=1,2,3$，$N=200\to400$ 的局部斜率 $-0.50$、$-1.00$、$-1.50$）；二維自迴避 $\theta=\frac{11}{24}$（$s=\frac23$）。
\item 動能項不變。\quad 答：(8.22) 的動能項寫成 $\frac3{2a'^2}$，$a'=a\,l_1/\sqrt{l_2}$；自迴避項乘 $\frac{l_1^d}{l_2^2}$。取 $l_1=\sqrt{l_2}$（$a'=a$），自迴避項乘 $l_2^{d/2-2}=l_2^{(d-4)/2}$。
\item 縮到一步。\quad 答：$l_2=\frac1N$ 時 $N'=1$、$a'=a$，$R'\sim a$，放大回去 $R=R'/l_1=a\sqrt N$；同時自迴避被乘上 $N^{(4-d)/2}$（$N=100$：1000、100、10、1、0.1）。
$d<4$ 時它隨 $N$ 變大，$\sqrt N$ 不可能是答案；$d>4$ 時它趨於 0，普通漫步的 $\sqrt N$ 是對的：上臨界維度 4。
\item 冪次型交互作用。\quad 答：$v_0|\vec r(s_1)-\vec r(s_2)|^{-q}$ 給 $\frac{l_1^q}{l_2^2}$，$l_1=\sqrt{l_2}$ 時 $l_2^{(q-4)/2}$，與 $d$ 無關：$q$ 取代了 $\delta$ 函數的 $d$。
$q<d$ 時是長程的，$q<4$ 就在任何維度都重要；$q>d$ 時短距離部分主導，等於接觸作用，又回到 $d_c=4$。合起來：重要 $\Leftrightarrow\min(d,q)<4$；Flory 版 $\nu=\frac3{q+2}$（$q<d$）。
\end{enumerate}}

% =====================================================================
\subsection{第 8 章的回顧與筆誤一覽}
\begin{center}
\small\setlength{\tabcolsep}{4pt}
\begin{tabular}{@{}llll@{}}
\toprule
主題 & 結果 & 書上的式子 & 動畫\\
\midrule
長尾巴 & 寬度 $\propto N^{1/(2-\delta)}$（1000 步：1822 對 28.06）、尾巴 $r^{-5+\delta}$ & (8.1)--(8.7) & W801\\
標度形式 & 塌縮 $(z_c-z)g\to\frac2{4+x^2}$；$\xi\propto(z_c-z)^{-\nu}$；$G=r^{-d+2-\eta}F(r(z_c-z)^\nu)$ & (8.8)--(8.14) & W802\\
最陡下降 & $C\sim e^{-(r/N^\nu)^{1/(1-\nu)}}$，Fisher 2 與 4；$\gamma=(2-\eta)\nu=\frac{43}{32}$ & (8.15)--(8.20) & W803\\
重新標度 & 自迴避項 $\times l_2^{(d-4)/2}$，上臨界維度 4；$\nu=\frac3{d+2}$（＝第 7 章 Flory） & (8.21)--(8.25) & W804\\
\bottomrule
\end{tabular}
\end{center}
\textbf{筆誤一覽}（給 2004 年版；〔注〕見各節）：下表在文中都用紅框標出，每一條都用程式核對過。
\begin{center}
\footnotesize\renewcommand{\arraystretch}{0.9}
\EightTabErrata
\end{center}
